
\section{辛同调}
\label{sec:symplectic-homology}

\hzsync{theorem}{14}%
\hzsync{equation}{159}%

我们在第 3 章中以\trDynamics{}的方式、借助\trHamiltonianSystem{}引入了特殊\trSymplectic%
\trCapacity{}$c_{0}(M,\omega)$。其存在性证明基于\trActionFunctional{}的一个特定%
\trCriticalPoint。这个\trCriticalPoint{}代表一条特定的\trPeriodicOrbit，于是问题随之%
而来：全体\trCriticalPoint{}连同其诸\trConnectingOrbit{}的结构是否会导致一个更为精细的%
\trSymplecticInvariant——它不再必然是一个数。的确，把\trCapacity{}$c_{0}$的构造与%
\trLoopSpace{}中的\trFloer{}同调构造相结合，H.~\trHofer{}与 A.~\trFloer{}在 [90] 中建立%
了一种\trSymplectic\trHomology{}理论。我们的目标是不加证明地勾画这一尚未定型的理论。

我们先回顾上一节的框架，为简单起见，考虑满足 $\pi_{2}(M)=0$ 的一个%
\trSymplecticManifold{} $(M,\omega)$。以 $\Omega(M)$ 表示参数化的可缩\trLoop{}%
$x:S^{1}\to M$ 之集。若 $H:S^{1}\times M\to\R$ 是一个时间\trPeriodic{}的%
\trHamiltonianFunction，则\trActionFunctional{} $\varphi_{H}:\Omega(M)\to\R$ 定义为
\[
\varphi_{H}(x)\ =\ -\int_{D}\bar{x}^{\,*}\omega\ +\ \int_{0}^{1}H\Bigl(t,x(t)\Bigr)\,dt\,.
\]
此处 $\bar{x}:D\to M$ 是 $x$ 到闭单位\trDisc{} $D$ 的一个延拓，满足%
$\bar{x}\bigl|_{\partial D}=x$。$\varphi_{H}$ 的诸\trCriticalPoint{}就是 $M$ 上%
\trHamiltonianVectorField{} $\dot{x}=X_{H}(x)$ 的可缩 $1$-\trPeriodicSolution。我们将假设它%
们全都\trNondegenerate，并把这个有限集记作 $\mathcal{P}_{H}$。以 $C(M,H)$ 表示由%
$\mathcal{P}_{H}$ 生成的 $\Z_{2}$ 上向量空间，即\footnote{原书印作「$C=\bigoplus C_{k}$」，未标明求和下标，与上一节 \eqref{eq:floer-chain-C} 之定义不一致，此处补 $k\in\Z$ 径改。——译者注}
\[
C\ =\ \bigoplus_{k\in\Z}C_{k}\,,\qquad C_{k}(M,H)\ =\ \operatorname{span}_{\Z_{2}}\Bigl\{\,x\in\mathcal{P}_{H}\ \Bigm|\ \mu(H,x)=k\,\Bigr\}\,.
\]
分次由内在地伴随诸\trNondegenerate\trPeriodicSolution{} $x\in\mathcal{P}_{H}$ 的指数%
$\mu(H,x)\in\Z$ 确定。对一个与 $\omega$ 相配的通有\trAlmostComplexStructure{} $J$，存在%
\trBoundaryOperator{} $\partial_{k}=\partial_{k}(M,H,J):C_{k}\to C_{k-1}$。它的定义方式是：%
按（模 $2$）计数连接指数差 $\mu(H,x^{-})-\mu(H,x^{+})=1$ 的两个\trPeriodicSolution{}的诸%
\trConnectingOrbit{} $\mathcal{M}(x^{-},x^{+},H,J)$。序对 $(C,\partial)$ 构成一个\trChainComplex，%
而我们把
\[
F_{*}(M,H,J)\ =\ \frac{\ker\partial}{\operatorname{im}\partial}
\]
称为通有对 $(H,J)$ 的\trFloer{}同调。现在转向 $(M,\omega)$ 的\trSymplectic%
\trHomology{}的构造，我们先作一个观察：对一个取定的通有对 $(H,J)$，\trFloer{}%
复形上还存在着一种由泛函 $\varphi_{H}$ 在 $\Omega(M)$ 上的值定义的附加%
\textbf{\trFiltration}（filtration）。我们对 $a\in\R\cup\{\infty\}$ 定义带分次的自由群
\[
\begin{aligned}
C^{a}(M,H)\ &=\ \bigoplus_{k\in\Z}C_{k}^{a}\\
C_{k}^{a}\ &=\ \operatorname{span}_{\Z_{2}}\Bigl\{\,x\in\mathcal{P}_{H}\ \Bigm|\ \mu(H,x)=k\ \text{且}\ \varphi_{H}(x)<a\,\Bigr\}\,.
\end{aligned}
\]
现在设 $u$ 是 $\mathcal{M}(x^{-},x^{+},H,J)$ 中的一条\trConnectingOrbit，并以%
$u_{s}\in\Omega$ 记诸\trLoop{} $u_{s}(t)=u(s,t)$。由变分结构，%
$\varphi_{H}(x^{+})\leq\varphi_{H}(u_{s})\leq\varphi_{H}(x^{-})$，于是位于两个不同水平之间的%
诸\trPeriodicOrbit{}连同其诸\trConnectingOrbit{}便构成\trFloer{}复形的一个子复形。更精确%
地说，$\partial=\partial(H,J)$ 满足
\[
\partial\ :\ C_{k}^{a}\ \longrightarrow\ C_{k-1}^{a}\,.
\]
因此，对 $-\infty<a\leq b\leq\infty$，用
\begin{equation}\label{eq:quotient-Cab}
C^{[a,b)}(H,J)\ =\ \frac{C^{b}(H,J)}{C^{a}(H,J)}\,,
\end{equation}
来定义商群，则边缘映射便诱导出一个\trBoundaryOperator{} $\partial_{k}:C_{k}^{[a,b)}\to C_{k-1}^{[a,b)}$，%
而我们可以定义
\begin{equation}\label{eq:relative-floer-group}
F_{k}^{[a,b)}(H,J)\ =\ \frac{\ker\bigl(\partial_{k}\bigr)}{\operatorname{im}\bigl(\partial_{k+1}\bigr)}\,.
\end{equation}
然而，这一\trFloer{}同调群依赖于通有对 $(H,J)$ 的选取，并不是 $M$ 的一个%
\trInvariant。我们的目的是构造一个由非常特殊的通有对组成的\textbf{\trDirectedSet}%
（directed set），使得相应群的\textbf{\trDirectedLimit}（directed limit）定义 $M$ 的一个%
\trSymplecticInvariant。通有对中容许\trHamiltonian{}函数的类，其原型便是第 3 章中构造%
$c_{0}(M,\omega)$ 时所用的那个类。

我们对可能带有\trBoundary{} $\partial M$ 的\trCompact\trSymplecticManifold{} $(M,\omega)$ 勾画这%
一构造。当 $\partial M\neq\varnothing$ 时，我们将假设 $\partial M$ 具有上文所定义的%
\trContactType。等价地，这要求存在 $\partial M$ 上的一个\trOneForm{} $\lambda$，它在%
$\partial M$ 上满足 $d\lambda=\omega$，并且对一切非零的 $\xi\in T_{x}(\partial M)$——它们%
对每个 $\eta\in T_{x}(\partial M)$ 满足 $\omega(\xi,\eta)=0$——都有 $\lambda(x,\xi)\neq0$。%
此外，与上文一样，为简单起见我们假设 $\pi_{2}(M)=0$。以 $\mathcal{H}$ 表示全体满足
\begin{equation}\label{eq:admissible-H}
\operatorname{supp}(H)\ \text{在 }S^{1}\times\bigl(M\setminus\partial M\bigr)\ \text{中为\trCompact{}的，且}\ H\bigl(S^{1}\times M\bigr)\subset(-\infty,0]\,.
\end{equation}
的\trSmooth{}时间\trPeriodic\trHamiltonianFunction{} $H:S^{1}\times M\to\R$ 组成的集合。%
显然，这个条件蕴含：$M$ 上 $X_{H}$ 的诸\trPeriodicSolution{}中有许多是%
\trDegenerate{}的且满足 $\varphi_{H}(x)=0$。因此，除 \eqref{eq:admissible-H} 之外，我们还%
要求 $X_{H}$ 的一切满足 $\varphi_{H}(x)<0$ 的可缩\trPeriodicSolution{}都\trNondegenerate。%
于是，若 $-\infty\leq a\leq b\leq0$，我们便可像上文那样定义相对\trFloer{}群%
$F^{[a,b)}(H,J)$。然而，此处并非每个 $J$ 都是容许的。

$J$ 的选取必须使得 $\mathcal{M}(x^{-},x^{+},H,J)$ 中的每个解 $u:\R\times S^{1}\to M$%
都远离\trBoundary：即当 $x^{-}$ 与 $x^{+}$ 均含于 $M\setminus\partial M$ 时，%
$u\bigl(\R\times S^{1}\bigr)\subset M\setminus\partial M$。这样的容许 $J$ 确实存在；此时就%
称\trBoundary{} $\partial M$ 为 D.~\trMcDuff{} [154] 意义下的\textbf{$J$-\trPseudoConvex}%
（$J$-pseudo convex）。我们将把满足上述条件的对 $(H,J)$ 称为容许对，并在容许对之%
集中引入一个\trPartialOrder，办法是规定（这并不是一个笔误）
\[
\bigl(H,J\bigr)\ \leq\ \bigl(K,\tilde{J}\bigr)\qquad\Longleftrightarrow\qquad H(t,x)\ \geq\ K(t,x)
\]
对一切 $(t,x)\in S^{1}\times M$ 成立。

对这样的对可以定义一个\textbf{\trMonotone\trHomotopy}（monotone homotopy）。它由一个对%
$(L,\hat{J})$ 组成，此处 $L:\R\times S^{1}\times M\to\R$，而\trAlmostComplexStructure{}%
$\hat{J}(s,t,x)$ 与\trSymplecticForm{} $\omega$ 相适配。此外，当 $s\to-\infty$ 时 $L=H$、%
$\hat{J}=J$，而当 $s\to+\infty$ 时 $L=K$、$\hat{J}=\tilde{J}$，并且%
$\partial L/\partial s\leq0$。给定 $(H,J)\leq(K,\tilde{J})$ 的一个单调同伦，便考虑%
$u:\R\times S^{1}\to M$ 的偏微分方程组
\[
\begin{aligned}
&\ \frac{\partial u}{\partial s}\ +\ \hat{J}(s,t,u)\,\frac{\partial u}{\partial t}\ +\ \bigl(\nabla_{\hat{J}}L\bigr)(s,t,u)\ =\ 0\\
&\ u(s,\cdot)\ \longrightarrow\ x\in\mathcal{P}_{H}\qquad(s\to-\infty)\\
&\ u(s,\cdot)\ \longrightarrow\ y\in\mathcal{P}_{K}\qquad(s\to+\infty)\,.
\end{aligned}
\]
利用单调性假设我们得出 $\varphi_{K}(y)\leq\varphi_{H}(x)$。这蕴含\trBoundaryOperator{}%
$\partial$ 保持实\trFiltration；这正是引入单调性假设的缘由。仿照\trFloer{}的构造%
去做，人们用解之集的组合学来构造 $\Z_{2}$ 向量空间之间的一个自然同态
\[
F^{[a,b)}(H,J)\ \longrightarrow\ F^{[a,b)}\bigl(K,\tilde{J}\bigr)
\]
它与所选的通有单调同伦无关。

这个偏序 $\leq$ 使容许对 $(H,J)$ 之集成为一个\trDirectedSet：因为给定 $(H_{1},J_{1})$ 与%
$(H_{2},J_{2})$，总存在一个容许对 $(H,J)$ 满足
\[
\bigl(H_{j},J_{j}\bigr)\ \leq\ \bigl(H,J\bigr)\qquad\text{对 }j=1,2\,.
\]
于是我们可以取极限，从而得到\trSymplectic\trHomology{}群
\[
S^{M}_{[a,b)}\ =\ \operatorname{dir.\,lim}\ F^{[a,b)}(H,J)\,.
\]
此外，若 $U\subset M$ 是一个开子集，且 $\mathcal{H}_{U}$ 由那些受「%
$\operatorname{supp}(H)\subset S^{1}\times(U\setminus\partial M)$」这一要求限制的容许对%
$(H,J)$ 组成，则我们可以在 $\mathcal{H}_{U}$ 上取\trDirectedLimit，得到 $U$ 的一个%
\trSymplectic\trHomology{}群，定义为
\[
S^{M}_{[a,b)}(U)\ :=\ \operatorname*{dir.\,lim}_{\mathcal{H}_{U}}\ F^{[a,b)}(H,J)\,.
\]
下面我们列举这些群迄今已知的一些性质。可以证明，存在 $\Z_{2}$ 向量空间之间的%
自然同态
\[
S^{M}_{[a,b)}(U)\ \longrightarrow\ S^{M}_{[a,b)}\,.
\]
此外，\textbf{短\trExactSequence}（short exact sequence）
\[
0\ \longrightarrow\ F^{[a,b)}(H,J)\ \longrightarrow\ F^{[a,c)}(H,J)\ \longrightarrow\ F^{[b,c)}(H,J)\ \longrightarrow\ 0
\]
（$-\infty\leq a\leq b\leq c\leq0$）给出了\trSymplectic\trHomology{}群的一个%
\textbf{\trExact{}三角形}（exact triangle）$\vartriangle_{a,b,c}(U)$：

\begin{center}
\begin{tikzpicture}[>=stealth,scale=1.0]
  \node (A) at (0,0) {$S^{[a,b)}(U)$};
  \node (B) at (4.6,0) {$S^{[a,c)}(U)$};
  \node (C) at (2.3,-1.8) {$S^{[b,c)}(U)$。};
  \draw[->] (A) -- (B);
  \draw[->] (B) -- (C);
  \draw[->] (C) -- (A);
  \node at (0.62,-1.22) {$\partial$};%
\end{tikzpicture}
\end{center}

给定一个\trSymplecticDiffeomorphism{} $\psi:M\to M$，对任意使 $H\in\mathcal{H}_{U}$ 的容许%
对 $(H,J)$，我们可以如下定义一个新的容许对 $(H_{\psi},J_{\psi})$：
\[
\begin{aligned}
H_{\psi}(t,x)\ &=\ H\Bigl(t,\psi^{-1}(x)\Bigr)\\
J_{\psi}(x)\ &=\ T\psi\Bigl(\psi^{-1}(x)\Bigr)\,J\Bigl(\psi^{-1}(x)\Bigr)\,T\psi^{-1}(x)\,.
\end{aligned}
\]
用它可证 $\psi$ 诱导一个自然同构
\[
\psi_{\#}\ :\ S^{M}_{[a,b)}(U)\ \xrightarrow{\ \cong\ }\ S^{M}_{[a,b)}\Bigl(\psi(U)\Bigr)\,.
\]
若 $U\subset V$，则\trInclusionMapping{}诱导一个同态
\[
\sigma_{V,U}\ :\ S^{M}_{[a,b)}(U)\ \longrightarrow\ S^{M}_{[a,b)}(V)\,.
\]
若 $\psi:M\to M$ 是满足 $\psi(U)\subset V$ 的一个\trSymplecticDiffeomorphism，我们定义诱%
导映射 $\psi_{*}$ 为
\[
\psi_{*}\ =\ \sigma_{V,\psi(U)}\circ\psi_{\#}\ :\ S^{M}_{[a,b)}(U)\ \longrightarrow\ S^{M}_{[a,b)}(V)\,.
\]
于是，若 $\varphi:M\to M$ 是另一个满足 $\varphi(V)\subset W$ 的%
\trSymplecticDiffeomorphism，则可证
\[
\begin{aligned}
\bigl(\varphi\circ\psi\bigr)_{*}\ &=\ \sigma_{W,\varphi\circ\psi(U)}\circ\bigl(\varphi\circ\psi\bigr)_{\#}\\
&=\ \Bigl(\sigma_{W,\varphi(V)}\circ\varphi_{\#}\Bigr)\circ\Bigl(\sigma_{V,\psi(U)}\circ\psi_{\#}\Bigr)\\
&=\ \varphi_{*}\circ\psi_{*}\,.
\end{aligned}
\]
此外，$M$ 的\trIdentityMap{}诱导恒同映射 $\operatorname{id}:S^{M}_{[a,b)}(U)\to S^{M}_{[a,b)}(U)$。%
最有用的性质是可如下陈述的\textbf{\trIsotopy\trInvariance}（isotopy invariance）：若%
$\psi_{s}:M\to M$（$0\leq s\leq1$）是\trSymplecticDiffeomorphism{}的一个\trSmooth{}族，%
满足 $\psi_{s}(U)\subset V$ 对一切 $s$ 成立，则 $(\psi_{s})_{*}:S^{M}_{[a,b)}(U)\to S^{M}_{[a,b)}(V)$%
不依赖于 $s$。

要计算上文勾画的诸\trSymplectic\trHomology{}群相当困难；有关计算与应用，我们%
参考 [50]。这些同调群与位于某个\trDomain{}的\trBoundary{}上或其附近的特定%
\trPeriodicOrbit{}相关联，这已不再令人惊讶。下述定理对此作了例证。

我们考虑一个带\trBoundary{} $\partial M\neq\varnothing$ 的\trCompact\trSymplecticManifold{}%
$(M,\omega)$，它是\textbf{\trExact\trSymplectic}（exact symplectic）的，即 $\omega=d\lambda$%
（$\lambda$ 为一个\trOneForm）。我们假设\trBoundary{}关于 $\lambda$ 具有切触型。%
于是我们可以如下定义\trActionSpectrum{} $\sigma(\partial M)\subset\R$\footnote{原书此定义之集合括号跨行裹住行文且句末多一枚闭括号，此处依下文 $d_{M}(\alpha)$ 同款记法径改为集合构造式。——译者注}：
\[
\sigma(\partial M)\ =\ \Bigl\{\,\Bigl|\int_{x}\lambda\Bigr|\ \Bigm|\ x\ \text{为}\ \partial M\ \text{上\trReeb{}向量场的周期解}\,\Bigr\}\,.
\]
注意迭代解也在考虑之列；特别地，对每个 $k\in\N$ 有 $k\cdot\sigma(\partial M)\subset\sigma(\partial M)$。

\begin{theorem}[K.~\trCieliebak、A.~\trFloer、H.~\trHofer{}与 K.~\trWysocki{ [52]}]\label{thm:action-spectrum-equal}
考虑两个正合\trSymplectic{}的\trCompact\trManifold{} $(M,d\lambda)$ 与 $(N,d\tau)$，其%
\trBoundary{} $\partial M$ 与 $\partial N$ 分别关于 $\lambda$ 与 $\tau$ 具有切触型。%
假设存在一个正合\trSymplecticDiffeomorphism
\[
\varphi\ :\ \mathring{M}\ \longrightarrow\ \mathring{N}
\]
（即 $\varphi^{*}\tau-\lambda=dF$，此处 $F:\mathring{M}\to\R$ 是一个\trSmooth{}函数）。%
若 $\partial M$ 与 $\partial N$ 上的诸\trPeriodicSolution{}连同其全部迭代都%
\trNondegenerate，则
\[
\sigma(\partial M)\ =\ \sigma(\partial N)\,.
\]
\end{theorem}

Ya.~\trEliashberg{}与 H.~\trHofer{}的例子 [75] 表明：在该定理的假设之下，$\partial M$ 与%
$\partial N$ 上的诸\trFlow{}未必是\trConjugate{}的，从而 $\varphi$ 不总能延拓成一个%
\trSymplecticDiffeomorphism{} $M\to N$。

\cref{thm:action-spectrum-equal}的证明基于下述非平凡的事实 [52]。设 $\alpha<0$，并%
定义
\[
d_{M}(\alpha)\ =\ \frac{1}{2}\lim_{\varepsilon\to0}\ \dim\,S^{\,[\alpha-\varepsilon,\alpha+\varepsilon)}_{M}\bigl(\mathring{M}\bigr)\,.
\]
则当 $-\alpha\notin\sigma(\partial M)$ 时 $d_{M}(\alpha)=0$；而
\[
d_{M}(\alpha)\ =\ \#\ \Bigl\{\ \text{$\partial M$ 上使\trAction{}}\ \Bigl|\int_{x}\lambda\Bigr|\ =\ -\alpha\ \text{的周期解}\ x\ \Bigr\}\,.
\]
由于对一切 $a<b$ 有 $S^{[a,b)}_{M}\bigl(\mathring{M}\bigr)=S^{[a,b)}_{N}\bigl(\mathring{N}\bigr)$，%
我们得出 $d_{M}(\alpha)=d_{N}(\alpha)$，定理便随之得证。有关证明的繁难细节以及这一%
结果的应用，我们参考即将问世的论文 [52]。

通过计算这些同调群，人们还证明了\trGromov{}关于\textbf{开\trPolydisc}%
（open polydiscs）之辛分类的一个猜想。对 $r=(r_{1},r_{2},\ldots,r_{n})$——其中%
$0<r_{1}\leq r_{2}\leq\cdots\leq r_{n}$——定义\trPolydisc{} $D(r)=B^{2}(r_{1})\times\cdots\times B^{2}(r_{n})$。

\begin{theorem}\label{thm:polydisc-symplectomorphic}
假设 $r$ 与 $r'$ 都是有序的，则开\trPolydisc{} $D(r)$ 与 $D(r')$ 彼此%
\trSymplectomorphism{}当且仅当 $r=r'$。
\end{theorem}

回顾一下，我们已在第 2 章中对特殊情形 $n=2$ 证明了这个陈述。该证明基于两个%
\trSymplecticInvariant：\trVolume{}与一个\trCapacity。$n\geq2$ 的证明连同其他应用见 [52]。%
上述\trSymplectic\trHomology{}构造的细节已由 K.~\trCieliebak、A.~\trFloer{}与%
H.~\trHofer{}在 [51] 中完成，文中还勾画了这一理论的一些可能推广。这个证明基于诸%
\trSymplecticManifold{}的\trLoopSpace{}中的\trAction{}原理，它把第 3 章的动力学%
\trCapacity{}的构造与\trFloer{}同调的构造结合起来，并附加一种实\trFiltration。%
对 $(\R^{2n},\omega_{0})$ 的开子集，L.~\trTraynor{} [212] 给出了另一种更晚近的%
\trSymplectic\trHomology{}理论构造；该构造基于有限维框架下的\trGeneratingFunction{}方法，后者由%
C.~\trViterbo{}在 [218] 中引入。

\textbf{结束语。}\quad 正如引言中已经指出的，我们只选取了\trSymplecticGeometry{}中%
少数几个相当初等的课题并作了详细解释。因此，我们愿意在此提及一些补充文献。%
对 $(\R^{2n},\omega_{0})$ 中的子集，I.~\trEkeland{}与 H.~\trHofer{}在 [68] 中于定义在%
整个 $\R^{2n}$ 上的\trSymplecticDiffeomorphism{}框架内构造了整整一列相对辛\trCapacity。%
这些\trCapacity{}具有\trOuterRegularity，与之相对照，第 3 章中内在定义的动力学%
\trCapacity{}则具有\trInnerRegularity。就在不久之前，S.~\trKuksin{}在 [130] 中把%
\trCapacity{}概念推广到由\trPartialDifferentialEquation{}所定义的某些特殊的无穷维%
\trHamiltonianSystem，例如 $S^{1}$ 上的非线性弦方程、$T^{n}$ 上的某些非线性\trSchrodinger%
方程。他特别证明了这类系统的诸\trFlow{}的一个非挤压结果。该\trCapacity{}的构造%
基于\trGalerkin{}逼近，其中使用了第 3 章中引入的有限维\trCapacity。\trKuksin{}的结果又%
促成了一个关于 $S^{1}$ 上非常特殊的非线性\trSchrodinger{}方程的非挤压结果，该方程的诸流%
映射并不是\trSymplectic\trHilbert{}空间中线性映射的\trCompact{}扰动，见 J.~\trBourgain{} [31]。%
把嵌入\trCapacity{}——第 2 章中称之为\trGromovWidth{}——加以推广，可以对每个%
整数 $k$ 如下定义一个新的\trCapacity：取 $k$ 个等\trBall{} $B(r)$ 的%
\trSymplecticEmbedding，使它们在 $(M,\omega)$ 中的像\trDisjoint，并定出半径 $r$ 的%
一个上界。这就引出与代数几何相关的\textbf{\trSymplectic\trPacking}（symplectic packing）问题，D.~\trMcDuff{}与%
L.~\trPolterovich{}在 [155] 中研究了它；至于有趣的\trPacking{}构造，我们参考%
L.~\trTraynor{} [213]。第 5 章中所描述的、$(\R^{2n},\omega_{0})$ 的紧支撑%
\trHamiltonianMap{}群上的内在 $L^{\infty}$-几何，已被 F.~\trLalonde{}与 D.~\trMcDuff{}%
在 [134] 与 [135] 中推广到一般的\trSymplecticManifold。结果表明，这一几何的基石，%
即能量-容量不等式，等价于非挤压定理。\trHofer{}度量下测地线上的共轭点，其间已%
由 I.~\trUstilovsky{}在 [214] 中作了研究。关于这些有趣进展的综述，我们参考%
F.~\trLalonde{} [133]。

正如已经强调过的，关于\trHamiltonianMap{}之\trFixedPoint{}的 V.I.~\trArnold{}猜想，%
对 $\dim M\geq8$ 的一般\trCompact\trSymplecticManifold{}仍然是开放的。有关\trFloer{}对这一%
问题的贡献的综述，我们参考 D.~\trMcDuff{} [153] 以及\trFloer{}纪念文集 [120]。%
晚近，\trLeHongVan{}与 K.~\trOno{} [137] 研究了\trArnold{}猜想的一个类似物，它针对的是%
虽未必是\trHamiltonianMap{}、但可通过一族\trSymplecticDiffeomorphism{}%
\trIsotopy{}于\trIdentityMap{}的\trSymplecticMapping{} $\varphi$。他们证明：对满足 $c_{1}\bigl|_{\pi_{2}(M)}=\lambda\,\omega\bigl|_{\pi_{2}(M)}$%
（$\lambda\geq0$）的\trCompact\trManifold，假设全部\trFixedPoint{}皆\trNondegenerate，%
则 $\varphi$ 的\trFixedPoint{}数目至少等于 $\Z_{2}$ 上由 $\varphi$ 的\trCalabi{}%
不变量 $\alpha$ 相伴的\textbf{\trNovikov\trHomology}（Novikov homology）%
$\operatorname{Nov}(\alpha)$ 的诸\trBettiNumber{}之和。在第 6 章中所考虑的\trTorus{}%
$T^{2n}$ 这一特殊情形，$\alpha\neq0$ 时 $\operatorname{Nov}(\alpha)=0$，故对\trTorus{}%
而言得不到任何新结果。

最后，我们的书中未论及的\trSymplecticGeometry{}中还有大量激动人心的进展与发现，%
例如伪全纯球面、填充技术、Lagrangian 交、切触几何、\trSymplecticManifold{}的例子%
等等；欲知其详，我们请读者参阅即将出版的下述诸书：D.~\trMcDuff{}与%
D.~\trSalamon{} [156]；M.~\trAudin{}、J.~\trLafontaine{} [13]；B.~\trAebischer{}、M.~\trBorer{}、%
Ch.~\trLeuenberger{}、M.~\trKalin{}、H.M.~\trReimann{} [3]；以及 C.~\trAbbas{}、H.~\trHofer{} [1]。%
